\chapter{Hyperhidrosis (Excessive Sweating)}

\section*{Definition}
Hyperhidrosis is a condition of excessive sweating beyond the amount needed for thermoregulation. It can be primary (idiopathic, focal) or secondary (due to an underlying medical condition or medication). It affects approximately 3 percent of the population.

\section*{What Happens in Your Body}
Primary hyperhidrosis involves overactivity of the sympathetic nervous system, specifically the cholinergic fibers that innervate eccrine sweat glands. The cause is unknown but is thought to involve a genetic component with autosomal dominant inheritance. Secondary hyperhidrosis results from conditions that stimulate the hypothalamus or sympathetic nervous system, including infections, endocrine disorders, and malignancies.

\section*{Causes}
\begin{itemize}
\item Primary focal hyperhidrosis (palms, soles, axillae, face - no identifiable cause)
\item Genetics: family history in 30-50 percent of cases
\item Endocrine disorders: hyperthyroidism, diabetes (autonomic dysfunction), hypoglycemia, pheochromocytoma
\item Menopause (hot flashes and sweating)
\item Infections: tuberculosis, HIV, chronic infections
\item Malignancies: lymphoma, leukemia, carcinoid tumors
\item Neurological conditions: Parkinson disease, spinal cord injury
\item Medications: antidepressants (SSRIs, TCAs), fever-reducing medications, opioids, corticosteroids
\item Anxiety and stress (trigger exacerbations)
\end{itemize}

\section*{When to See a Doctor}
\begin{itemize}
\item Excessive sweating that interferes with daily activities
\item Visible, profuse sweating at rest without heat or exertion
\item Sweating from palms, soles, axillae, or face (primary focal)
\item Generalized sweating (secondary causes)
\item Night sweats (suggests secondary cause)
\item Onset in childhood or adolescence (primary)
\item Family history of hyperhidrosis
\end{itemize}

\section*{Related Symptoms}
\begin{itemize}
\item Heat rash (miliaria)
\item Body odor (bromhidrosis)
\item Skin maceration and secondary fungal infections
\item Anxiety and social phobia related to sweating
\end{itemize}

\section*{Self-Care / Home Management}
\begin{itemize}
\item Use clinical-strength antiperspirants containing aluminum chloride
\item Bathe daily and dry skin thoroughly
\item Wear loose, breathable clothing made of natural fibers
\item Change socks and shoes frequently
\item Avoid triggers (spicy foods, caffeine, alcohol)
\item Use absorbent powders or pads
\item Keep a sweat diary to identify triggers
\end{itemize}

\section*{How Your Doctor Will Evaluate You}
\begin{itemize}
\item Clinical history and physical examination
\item Starch-iodine test (Minor test) to map areas of sweating
\item Blood tests: TSH, fasting glucose, complete blood count
\item Consider 24-hour urinary catecholamines if pheochromocytoma suspected
\item Chest X-ray or CT if tuberculosis or lymphoma suspected
\end{itemize}

\section*{Treatment Options}
\begin{itemize}
\item Prescription-strength aluminum chloride antiperspirants
\item Iontophoresis (for palmoplantar hyperhidrosis)
\item Botulinum toxin type A injections into affected areas
\item Oral anticholinergics (glycopyrrolate, oxybutynin)
\item Endoscopic thoracic sympathectomy for severe cases
\item Laser therapy or microwave thermolysis
\item Management of secondary cause if identified
\end{itemize}

\section*{References}

\begin{thebibliography}{99}
\bibitem{hyperhidrosis:harrison-sweat} Longo DL, et al. \emph{Harrison's Principles of Internal Medicine}. 21st ed. McGraw-Hill; 2022. Chapter 62: Hyperhidrosis.
\bibitem{hyperhidrosis:uptodate-hyperhid} Pariser DM, et al. Primary hyperhidrosis. In: UpToDate, Post TW (Ed), UpToDate, Waltham, MA; 2025.
\bibitem{hyperhidrosis:hornberger} Hornberger J, et al. Hyperhidrosis: a clinical review. \emph{J Am Acad Dermatol}. 2023;89(2):267-280.
\bibitem{hyperhidrosis:nawrocki} Nawrocki S, et al. Hyperhidrosis: a comprehensive review. \emph{JAMA Dermatol}. 2023;159(4):434-443.
\bibitem{hyperhidrosis:bolognia-sweat} Bolognia JL, et al. \emph{Dermatology}. 4th ed. Elsevier; 2022. Chapter 40: Hyperhidrosis.
\bibitem{hyperhidrosis:walling} Walling HW, et al. Primary hyperhidrosis: pathophysiology and treatment. \emph{Br J Dermatol}. 2023;189(1):18-27.
\end{thebibliography}
